% GENERATED FILE. Edit canonical lessons, then run npm run generate:books.
% canonical-source-hash: fnv1a64:d5fee06364bf271c
% canonical-lessons: TE-C493-samagri, TE-C493-sagu, TE-C493-tippu, TE-C493-cediripovu, TE-C493-venudirugu, TE-R493-review-486-493, TE-R493-review-final-493

\chapter{Heading Home}
\label{ch:te-heading-home}
\hlchaptermodality{493}

\begin{chapteropening}
\textbf{By the end of this chapter:} \emph{I can say equipment, to go on, to turn round, to scatter and to turn back.}

Complete the last lesson of chapter 493: I can say equipment, to go on, to turn round, to scatter and to turn back.
\end{chapteropening}

\section[\texorpdfstring{\te{సామగ్రి} (sāmagri)}{సామగ్రి (sāmagri)}]{\te{సామగ్రి} (sāmagri) — equipment; materials}

\label{lesson:TE-C493-samagri}



\begin{quote}
Before the new one: say the Telugu for a false step, then the Telugu for confusion.
\end{quote}

\subsection*{You'll want to know: \te{సామగ్రి}}
\textbf{\te{సామగ్రి}} — \emph{sāmagri} — \textquotedblleft{}equipment; materials\textquotedblright{}.

\begin{quote}
\te{కొండ} \te{ఎక్కడానికి} \te{కావలసిన} \te{సామగ్రి} \te{అంతా} \te{సర్దాం}.
\end{quote}

\begin{quote}
\emph{koṇḍa ekkaḍāniki kāvalasina sāmagri antā sardāṁ.}
\end{quote}

\begin{quote}
We packed all the gear we needed for the climb.
\end{quote}

Things needed for a task: \te{వంట} \te{సామగ్రి}, cooking things; \te{కట్టడ} \te{సామగ్రి}, building materials.

\begin{cousinweb}
\textbf{\te{సామగ్రి}} (\emph{sāmagri}) is borrowed from Sanskrit \emph{sāmagrī}, \textquotedblleft{}totality, the whole of what is needed\textquotedblright{} (Monier-Williams).
\end{cousinweb}

\begin{grammarlens}[title={What you've built}]
One more word for getting around.
\end{grammarlens}

\subsection*{Your turn}
\begin{itemize}
\raggedright
  \item \emph{Say it:} \emph{sāmagri}
  \item \emph{Say it:} \emph{sāmagri}, once more
  \item \emph{Say it:} \emph{koṇḍa ekkaḍāniki kāvalasina sāmagri antā sardāṁ.} — We packed all the gear we needed for the climb.
  \item \emph{Recall:} say the Telugu for a false step, then the Telugu for confusion, then say \emph{sāmagri} again
\end{itemize}

\begin{tcolorbox}[breakable,colback=teal!4,colframe=teal!35!black,title={Before you move on}]
What does \te{సామగ్రి} mean? (\textquotedblleft{}Equipment; materials\textquotedblright{}.) Say it once more.
\end{tcolorbox}
\section[\texorpdfstring{\te{సాగు} (sāgu)}{సాగు (sāgu)}]{\te{సాగు} (sāgu) — to go on; to proceed}

\label{lesson:TE-C493-sagu}



\begin{quote}
Before the new one: say the Telugu for confusion, then the Telugu for equipment.
\end{quote}

\subsection*{You'll want to know: \te{సాగు}}
\textbf{\te{సాగు}} — \emph{sāgu} — \textquotedblleft{}to go on; to proceed\textquotedblright{}.

\begin{quote}
\te{మా} \te{ప్రయాణం} \te{హాయిగా} \te{సాగింది}.
\end{quote}

\begin{quote}
\emph{mā prayāṇaṁ hāyigā sāgindi.}
\end{quote}

\begin{quote}
Our journey went on pleasantly.
\end{quote}

Past (I …): \te{సాగాను}, \emph{sāgānu}.

Said of a journey, a meeting or a life: \te{సాఫీగా} \te{సాగింది}, it went smoothly. As a noun \te{సాగు} is farming.

\begin{cousinweb}
\textbf{\te{సాగు}} (\emph{sāgu}) is native Dravidian: DEDR 2430 lists Telugu \emph{sāgu}, \textquotedblleft{}to go on, to proceed, to continue\textquotedblright{}.
\end{cousinweb}

\begin{grammarlens}[title={What you've built}]
One more word for getting around.
\end{grammarlens}

\subsection*{Your turn}
\begin{itemize}
\raggedright
  \item \emph{Say it:} \emph{sāgu}
  \item \emph{Say it:} \emph{sāgu}, once more
  \item \emph{Say it:} \emph{mā prayāṇaṁ hāyigā sāgindi.} — Our journey went on pleasantly.
  \item \emph{Recall:} say the Telugu for confusion, then the Telugu for equipment, then say \emph{sāgu} again
\end{itemize}

\begin{tcolorbox}[breakable,colback=teal!4,colframe=teal!35!black,title={Before you move on}]
What does \te{సాగు} mean? (\textquotedblleft{}To go on; to proceed\textquotedblright{}.) Say it once more.
\end{tcolorbox}
\section[\texorpdfstring{\te{తిప్పు} (tippu)}{తిప్పు (tippu)}]{\te{తిప్పు} (tippu) — to turn round; to twist}

\label{lesson:TE-C493-tippu}



\begin{quote}
Before the new one: say the Telugu for equipment, then the Telugu for to go on.
\end{quote}

\subsection*{You'll want to know: \te{తిప్పు}}
\textbf{\te{తిప్పు}} — \emph{tippu} — \textquotedblleft{}to turn round; to twist\textquotedblright{}.

\begin{quote}
\te{డ్రైవరు} \te{బండిని} \te{వెనక్కి} \te{తిప్పాడు}.
\end{quote}

\begin{quote}
\emph{ḍraivaru baṇḍini venakki tippāḍu.}
\end{quote}

\begin{quote}
The driver turned the cart back.
\end{quote}

Past (I …): \te{తిప్పాను}, \emph{tippānu}.

The causative of \te{తిరుగు}, to turn: to make something turn. \te{తాళంచెవి} \te{తిప్పు}, to turn a key.

\begin{cousinweb}
\textbf{\te{తిప్పు}} (\emph{tippu}) is native Dravidian, usually traced to the family of \te{తిరుగు}, \textquotedblleft{}to turn, to go round\textquotedblright{} (DEDR 3246), as its causative.
\end{cousinweb}

\begin{grammarlens}[title={What you've built}]
One more word for getting around.
\end{grammarlens}

\subsection*{Your turn}
\begin{itemize}
\raggedright
  \item \emph{Say it:} \emph{tippu}
  \item \emph{Say it:} \emph{tippu}, once more
  \item \emph{Say it:} \emph{ḍraivaru baṇḍini venakki tippāḍu.} — The driver turned the cart back.
  \item \emph{Recall:} say the Telugu for equipment, then the Telugu for to go on, then say \emph{tippu} again
\end{itemize}

\begin{tcolorbox}[breakable,colback=teal!4,colframe=teal!35!black,title={Before you move on}]
What does \te{తిప్పు} mean? (\textquotedblleft{}To turn round; to twist\textquotedblright{}.) Say it once more.
\end{tcolorbox}
\section[\texorpdfstring{\te{చెదిరిపోవు} (cediripōvu)}{చెదిరిపోవు (cediripōvu)}]{\te{చెదిరిపోవు} (cediripōvu) — to scatter; to disperse}

\label{lesson:TE-C493-cediripovu}



\begin{quote}
Before the new one: say the Telugu for to go on, then the Telugu for to turn round.
\end{quote}

\subsection*{You'll want to know: \te{చెదిరిపోవు}}
\textbf{\te{చెదిరిపోవు}} — \emph{cediripōvu} — \textquotedblleft{}to scatter; to disperse\textquotedblright{}.

\begin{quote}
\te{వర్షం} \te{రాగానే} \te{జనం} \te{చెదిరిపోయారు}.
\end{quote}

\begin{quote}
\emph{varṣaṁ rāgānē janaṁ cediripōyāru.}
\end{quote}

\begin{quote}
As soon as the rain came the crowd scattered.
\end{quote}

Past (it …): \te{చెదిరిపోయింది}, \emph{cediripōyindi}.

Said of a crowd, of clouds and of a dream. Its subject is a thing, so the past shown is the third person: \te{చెదిరిపోయింది}, it scattered.

\begin{cousinweb}
\textbf{\te{చెదిరిపోవు}} (\emph{cediripōvu}) is native Dravidian: DEDR 1546 lists Telugu \emph{cedaru}, \textquotedblleft{}to be scattered, dispersed\textquotedblright{}, here with \te{పోవు}, \textquotedblleft{}to go\textquotedblright{} (DEDR 4572).
\end{cousinweb}

\begin{grammarlens}[title={What you've built}]
One more word for getting around.
\end{grammarlens}

\subsection*{Your turn}
\begin{itemize}
\raggedright
  \item \emph{Say it:} \emph{cediripōvu}
  \item \emph{Say it:} \emph{cediripōvu}, once more
  \item \emph{Say it:} \emph{varṣaṁ rāgānē janaṁ cediripōyāru.} — As soon as the rain came the crowd scattered.
  \item \emph{Recall:} say the Telugu for to go on, then the Telugu for to turn round, then say \emph{cediripōvu} again
\end{itemize}

\begin{tcolorbox}[breakable,colback=teal!4,colframe=teal!35!black,title={Before you move on}]
What does \te{చెదిరిపోవు} mean? (\textquotedblleft{}To scatter; to disperse\textquotedblright{}.) Say it once more.
\end{tcolorbox}
\section[\texorpdfstring{\te{వెనుదిరుగు} (venudirugu)}{వెనుదిరుగు (venudirugu)}]{\te{వెనుదిరుగు} (venudirugu) — to turn back; to head back}

\label{lesson:TE-C493-venudirugu}



\begin{quote}
Before the new one: say the Telugu for to turn round, then the Telugu for to scatter.
\end{quote}

\subsection*{You'll want to know: \te{వెనుదిరుగు}}
\textbf{\te{వెనుదిరుగు}} — \emph{venudirugu} — \textquotedblleft{}to turn back; to head back\textquotedblright{}.

\begin{quote}
\te{వర్షం} \te{వల్ల} \te{మేము} \te{మధ్యలోనే} \te{వెనుదిరిగాం}.
\end{quote}

\begin{quote}
\emph{varṣaṁ valla mēmu madhyalōnē venudirigāṁ.}
\end{quote}

\begin{quote}
Because of the rain we turned back halfway.
\end{quote}

Past (I …): \te{వెనుదిరిగాను}, \emph{venudirigānu}.

\te{వెను}, back, + \te{తిరుగు}, to turn: to turn round and go back the way you came. Said of travellers, and of a side that gives up.

\begin{cousinweb}
\textbf{\te{వెనుదిరుగు}} (\emph{venudirugu}) is native Dravidian: \te{వెను}, \textquotedblleft{}the back\textquotedblright{}, which DEDR 5488 lists as Telugu \emph{ven(n)u} beside \emph{venuka}, \textquotedblleft{}behind\textquotedblright{}, and \te{తిరుగు}, \textquotedblleft{}to turn, to go round\textquotedblright{} (DEDR 3246), whose \te{త} becomes \te{ద} inside the compound.
\end{cousinweb}

\begin{grammarlens}[title={What you've built}]
One more word for getting around.
\end{grammarlens}

\subsection*{Your turn}
\begin{itemize}
\raggedright
  \item \emph{Say it:} \emph{venudirugu}
  \item \emph{Say it:} \emph{venudirugu}, once more
  \item \emph{Say it:} \emph{varṣaṁ valla mēmu madhyalōnē venudirigāṁ.} — Because of the rain we turned back halfway.
  \item \emph{Recall:} say the Telugu for to turn round, then the Telugu for to scatter, then say \emph{venudirugu} again
\end{itemize}

\begin{tcolorbox}[breakable,colback=teal!4,colframe=teal!35!black,title={Before you move on}]
What does \te{వెనుదిరుగు} mean? (\textquotedblleft{}To turn back; to head back\textquotedblright{}.) Say it once more.
\end{tcolorbox}
\section[(dialogue)]{Getting around, the last stretch}

\label{lesson:TE-R493-review-486-493}



\begin{quote}
Say the Telugu for a measurement and size.
\end{quote}

\subsection*{Your turn}
Say these aloud:

\begin{itemize}
\raggedright
  \item a measurement, size, a sample, value, to weigh
  \item to spend money, to save money, demand, to be available, to sell out
  \item wholesale, a stock, a storehouse, a shortage, an auction
  \item a pickpocket, adulteration, to be cheated, to rot, to be left over
  \item to spoil, to get torn, to get grubby, to split, to break off
  \item to wear down, to go mouldy, to fade, to crack, to deteriorate
  \item to be mistaken, to falter, to estimate, a false step, confusion
  \item equipment, to go on, to turn round, to scatter, to turn back
\end{itemize}

\begin{tcolorbox}[breakable,colback=teal!4,colframe=teal!35!black,title={Before you move on}]
A measurement? (\textbf{\te{కొలత}}, \emph{kolata}.) Size? (\textbf{\te{పరిమాణం}}, \emph{parimāṇaṁ}.) A sample? (\textbf{\te{నమూనా}}, \emph{namūnā}.) Value? (\textbf{\te{విలువ}}, \emph{viluva}.) To weigh? (\textbf{\te{తూచు}}, \emph{tūcu}.) To spend money? (\textbf{\te{వెచ్చించు}}, \emph{vecciñcu}.) To save money? (\textbf{\te{ఆదా} \te{చేయు}}, \emph{ādā cēyu}.) Demand? (\textbf{\te{గిరాకీ}}, \emph{girākī}.) To be available? (\textbf{\te{లభించు}}, \emph{labhiñcu}.) To sell out? (\textbf{\te{అమ్ముడుపోవు}}, \emph{ammuḍupōvu}.) Wholesale? (\textbf{\te{టోకు}}, \emph{ṭōku}.) A stock? (\textbf{\te{నిల్వ}}, \emph{nilva}.) A storehouse? (\textbf{\te{గిడ్డంగి}}, \emph{giḍḍaṅgi}.) A shortage? (\textbf{\te{కొరత}}, \emph{korata}.) An auction? (\textbf{\te{వేలం}}, \emph{vēlaṁ}.) A pickpocket? (\textbf{\te{జేబుదొంగ}}, \emph{jēbudoṅga}.) Adulteration? (\textbf{\te{కల్తీ}}, \emph{kaltī}.) To be cheated? (\textbf{\te{మోసపోవు}}, \emph{mōsapōvu}.) To rot? (\textbf{\te{కుళ్ళు}}, \emph{kuḷḷu}.) To be left over? (\textbf{\te{మిగులు}}, \emph{migulu}.) To spoil? (\textbf{\te{చెడగొట్టు}}, \emph{ceḍagoṭṭu}.) To get torn? (\textbf{\te{చిరిగిపోవు}}, \emph{cirigipōvu}.) To get grubby? (\textbf{\te{మాసిపోవు}}, \emph{māsipōvu}.) To split? (\textbf{\te{చీలిపోవు}}, \emph{cīlipōvu}.) To break off? (\textbf{\te{విరిగిపోవు}}, \emph{virigipōvu}.) To wear down? (\textbf{\te{అరిగిపోవు}}, \emph{arigipōvu}.) To go mouldy? (\textbf{\te{బూజుపట్టు}}, \emph{būjupaṭṭu}.) To fade? (\textbf{\te{వెలిసిపోవు}}, \emph{velisipōvu}.) To crack? (\textbf{\te{పగిలిపోవు}}, \emph{pagilipōvu}.) To deteriorate? (\textbf{\te{క్షీణించు}}, \emph{kṣīṇiñcu}.) To be mistaken? (\textbf{\te{పొరబడు}}, \emph{porabaḍu}.) To falter? (\textbf{\te{తడబడు}}, \emph{taḍabaḍu}.) To estimate? (\textbf{\te{అంచనా} \te{వేయు}}, \emph{añcanā vēyu}.) A false step? (\textbf{\te{తప్పటడుగు}}, \emph{tappaṭaḍugu}.) Confusion? (\textbf{\te{గందరగోళం}}, \emph{gandaragōḷaṁ}.) Equipment? (\textbf{\te{సామగ్రి}}, \emph{sāmagri}.) To go on? (\textbf{\te{సాగు}}, \emph{sāgu}.) To turn round? (\textbf{\te{తిప్పు}}, \emph{tippu}.) To scatter? (\textbf{\te{చెదిరిపోవు}}, \emph{cediripōvu}.) To turn back? (\textbf{\te{వెనుదిరుగు}}, \emph{venudirugu}.)
\end{tcolorbox}
\section[(dialogue)]{Heading Home — once more}

\label{lesson:TE-R493-review-final-493}



\begin{quote}
Say the Telugu for equipment and to go on.
\end{quote}

\subsection*{Your turn}
\begin{itemize}
\raggedright
  \item \emph{Say it:} equipment, to go on, to turn round, to scatter, to turn back
\end{itemize}

\begin{tcolorbox}[breakable,colback=teal!4,colframe=teal!35!black,title={Before you move on}]
Equipment? (\textbf{\te{సామగ్రి}}, \emph{sāmagri}.) To go on? (\textbf{\te{సాగు}}, \emph{sāgu}.) To turn round? (\textbf{\te{తిప్పు}}, \emph{tippu}.) To scatter? (\textbf{\te{చెదిరిపోవు}}, \emph{cediripōvu}.) To turn back? (\textbf{\te{వెనుదిరుగు}}, \emph{venudirugu}.)
\end{tcolorbox}
